% Revised 2026-08-28 after full-text verification of farea2025qcpinn and
% hqpinn2025flow. Key corrections vs. the earlier abstract-only draft:
% (1) farea2025qcpinn DOES report per-iteration timing (15-40x), just not
%     in the abstract, and its own results are already mixed (worse on 2/5
%     PDEs) -- don't overstate it as a uniformly positive prior result;
% (2) hqpinn2025flow's own measured max ratio is ~82x, not a clean 100x --
%     that rounded figure is their own summary of their own Table A3, while
%     a separate "two orders of magnitude" claim earlier in their paper is
%     attributed to OTHER authors, not this paper's own experiment;
% (3) "barren plateau" is NOT a term used in hqpinn2025flow at all (they
%     say "trainability" / "excessive expressibility" instead); the one
%     paper that does invoke barren plateaus explicitly is
%     farea2025qcpinn's continuous-variable circuit discussion -- which is
%     a DIFFERENT circuit paradigm from the qubit-based one we build on,
%     so don't cite that barren-plateau discussion as if it were about the
%     same architecture family.

\subsection{Prior positive results are narrower than they first appear}
None of \citet{farea2025qcpinn}, \citet{hqpinn2025flow}, or
\citet{qcpinn2025reservoir} is a uniform "quantum wins" result once read
past the abstract: all three contain within-paper negative or mixed cases
-- cavity flow and the wave equation for the first, non-harmonic and
transonic regimes for the second, no circuit topology winning across all
four benchmarks (plus an explicit disclaimer against causal claims) for
the third. We read our six-scenario result not as contradicting these
papers but as a more consistent, more uniform version of a pattern they
already show in part: where they find a mix depending on PDE, regime, or
topology, we find the quantum branch losing on every single one of six
matched scenarios on a problem family -- a diffusion equation under three
boundary-condition types, in both forward and inverse form -- that none of
the three tests at all: all three restrict to forward problems with
Dirichlet (and, in one case, additionally Neumann) conditions, and none
includes an inverse/parameter-estimation variant. It is also worth being
precise about how much independent weight \citet{qcpinn2025reservoir}
carries: it is not an arms-length third data point, since it explicitly
reuses \citeauthor{farea2025qcpinn}'s open-source implementation for a new
PDE domain. Its convergence with \citet{hqpinn2025flow} on the qualitative
claim that a quantum branch's advantage is regime- or topology-dependent
rather than uniform therefore reflects two independent author teams
agreeing on that shape, not three independent confirmations of any
specific number.

\subsection{Wall-clock accounting: sharper than the literature, but comparable}
\citet{farea2025qcpinn} report per-iteration time (15--40$\times$ slower
than classical, across all five of their benchmarks) and
\citet{hqpinn2025flow} report per-epoch time for one of their three
benchmarks (17--82$\times$, which they themselves round to "up to
100$\times$"). \citet{qcpinn2025reservoir} and \citet{hqcpinn2026hydro} sit
at the far end of this spectrum: the former runs on a fully specified CPU
(two Intel Xeon E5-2680 v4, 28 cores, 128GB RAM) and explicitly
acknowledges in prose that "quantum simulation has inherent overhead on
classical hardware," yet reports no timing measurement whatsoever; the
latter reports an epoch-count "speedup" (fewer epochs to reach a
validation-loss threshold) without ever mentioning wall-clock cost at all,
despite using parameter-shift-rule gradients -- known to require multiple
extra circuit evaluations per step relative to backprop -- which makes an
epoch-count comparison an especially weak proxy for real training cost in
that specific case. Two of four closely related papers reporting zero
timing data, and a third reporting it for only one of three benchmarks, is
itself worth stating as a field-wide pattern rather than a coincidence
across independent author teams. Our own 53--89$\times$ wall-clock ratio is
in the same direction as the two papers that do measure it, and a similar
order of magnitude -- not a qualitatively new finding -- but we measure it
as a primary outcome across all six of our scenarios rather than as a
secondary column reported for some subset of benchmarks, or omitted
entirely. Our ratio runs somewhat higher than Farea et al.'s 15--40$\times$;
one plausible source is circuit size, since their best-performing
configuration uses only 5 qubits and 1 layer, smaller than our 4-qubit,
2-layer baseline (and our own ablation sweeps up to 6 qubits and 4 layers,
which should be commensurately slower still). However, none of the three
papers states enough hardware detail to rule out a hardware difference as
a contributing factor instead or in addition -- Farea et al.\ report a
6148-CPU, 192GB cluster node; the flow paper reports only "distributed
CPU, no GPU"; the reservoir paper specifies its CPU but has no timing to
compare against at all -- so this is a limit on how precisely wall-clock
ratios can be compared \emph{across} papers, not only a feature of our own
results.

\subsection{A smooth problem that still loses -- and why "barren plateau" needs care}
\citet{hqpinn2025flow}'s central thesis is that harmonicity of the target
solution predicts whether the quantum branch helps, not problem difficulty
in general. Our diffusion PDE has closed-form sine/cosine solutions and so
sits on the "harmonic" side of their distinction, yet our quantum branch
loses in all six scenarios, including the smoothest ones. This is a
genuine point of tension worth naming directly: either the
harmonic/non-harmonic axis from their compressible-flow, Fourier-series
argument does not transfer to a diffusion PDE encoded via AngleEmbedding,
or some other factor in our setup -- data encoding, our ring-CNOT ansatz
versus their tunable CRX/CRZ entanglers (Section~\ref{sec:method}), or
optimizer choice -- is the actual confound; our results cannot distinguish
between these. We are also careful with terminology here, since neither
prior paper's account of trainability loss for the circuit family closest
to ours uses "barren plateau": \citet{hqpinn2025flow} calls it "excessive
expressibility," based on final-loss curves rather than a direct gradient
measurement, and \citet{farea2025qcpinn} uses "barren plateau" explicitly
but only for its continuous-variable (qumode) circuits, a different
paradigm from the qubit-based ansatz both they and we build on. Our own
phase-2 gradient-norm measurement (Section~\ref{sec:experiments}) is, as
far as this literature review found, the only \emph{direct} circuit-only
gradient measurement across a qubit-count/depth sweep for a qubit-based
hybrid PINN. We offer this as a contribution, but state plainly, as
already noted in Section~\ref{sec:experiments}, that our own numbers do
not cleanly resolve the underlying question either.

One theoretical result changes how that non-resolution should be framed,
though: \citet{holmes2022expressibility} show that gradient magnitude and
ansatz expressibility are not independent, competing explanations but two
faces of the same relationship -- more expressive ansätze induce flatter
cost landscapes and correspondingly smaller typical gradients. Read
through that lens, \texttt{wide} having both the best accuracy \emph{and}
the smallest final gradient norm among your four quantum configurations is
not the contradiction it looks like under a naive "expressivity good,
vanishing gradient bad" framing -- it is closer to what
\citet{holmes2022expressibility}'s bound would predict for a more
expressive circuit that has not yet entered a true (exponential,
system-size-dependent) barren plateau: still trainable, but with a flatter
landscape than \texttt{narrow} or \texttt{baseline}. We stop short of
calling this a confirmed barren plateau, however: our qubit-count range
(2--6) is almost certainly too small for the asymptotic, exponential-decay
regime that the barren-plateau literature
\citep{mcclean2018barren,holmes2022expressibility} actually characterizes.
What our sweep shows is consistent with the small-system-size end of the
expressibility-gradient relationship those results describe -- a flatter
but still-trainable landscape as the circuit widens -- not evidence for or
against an eventual barren plateau at larger qubit counts, which our range
cannot distinguish from ordinary flattening.

\subsection{Limitations}
\begin{itemize}
  \item All quantum-circuit simulation is exact and noiseless
    (\texttt{default.qubit}); results say nothing about performance on
    real NISQ hardware, where shot noise and decoherence would likely
    widen rather than close the gap. (Farea et al.'s feasibility analysis
    of their Cross-mesh circuit under a stated noise model is the closest
    precedent for what a hardware-realistic version of this comparison
    would need to account for.)
  \item Both branches were run on CPU. This asymmetry is worth reading
    carefully rather than at face value: the classical branch could in
    practice be GPU-accelerated, which would shrink its training time
    further, while the quantum branch's cost is dominated by classically
    simulating the circuit and would not benefit comparably from the same
    hardware -- so a GPU-enabled comparison would likely widen, not narrow,
    the wall-clock gap reported here. Under future fault-tolerant quantum
    hardware this asymmetry could of course reverse, but that is a
    statement about hardware that does not yet exist, not about the
    noiseless-simulator comparison in this paper. We also note that
    neither \citet{farea2025qcpinn} nor \citet{hqpinn2025flow} give full
    hardware specifications sufficient for a precise cross-paper timing
    comparison -- a field-wide reporting gap, not a limitation specific to
    our study.
  \item Single PDE family (linear diffusion, 3 BC types); results may not
    transfer to nonlinear or higher-dimensional PDEs, or to the
    discontinuous/shock-type solutions where \citet{hqpinn2025flow}
    independently find classical dominance for different reasons.
  \item The two ablations use 3 seeds per (scenario, configuration), but
    the headline classical-vs-quantum comparison in
    Table~\ref{tab:main-results} uses a single seed per scenario -- it is
    the only result in the paper reported as a point estimate rather than
    a mean $\pm$ std, and the largest-magnitude claims (the 53--89$\times$
    wall-clock ratios) come from that single-seed table. This is a real
    limitation rather than a disclosure formality: the ablations show
    run-to-run standard deviations on the same order as some of the
    effects being compared, so we cannot yet rule out that the main
    comparison's specific ratios would shift, within their same broad
    range, under a different seed. Multi-seed replication of
    Table~\ref{tab:main-results} is the clearest concrete next step before
    this result is submission-ready.
  \item Phase-2 ablation varies qubit count and depth but not entangler
    type (fixed CNOT throughout) or data-encoding scheme --
    \citet{farea2025qcpinn} suggest entangler tunability (CRX/CRZ vs.\
    fixed CNOT) affects convergence reliability in their own qubit-circuit
    sweep, which we do not test and flag as the clearest concrete next
    ablation.
\end{itemize}
